\section{Limitations and Conclusion}\label{sec:concl}

\subsection{Limitations and Threats to Validity}\label{sec:limits}
\textit{Internal validity.} Forty runs per workload give precise medians but
imprecise tail percentiles. Interleaving, warm-up and keeping outliers limit,
but do not remove, the influence of background activity. The planned groups of
15--20 teachers are small: a power of 0.80 is reached only for large differences
($d \gtrsim 0.9$--1.1, \autoref{tab:power}), so smaller ones may be missed.
\textit{External validity.} All measurements come from one machine and one
container, so absolute times describe this environment, and the respondents are
synthetic. The teachers are recruited by convenience sampling from several
schools of one city, so the results may not generalise to other cities or to
all schools, and those who agree to take part may differ from those who decline.
\textit{Construct validity.} Latency and throughput are proxies for how practical
the pipeline is. The MBI is a self-report questionnaire, and the index $B$
(\autoref{eq:index}) is relative to the sample, so it compares the two groups
but is not an absolute level of burnout.

\subsection{Conclusion and Future Work}\label{sec:future}
We presented the protocol of a planned comparison of teacher burnout in state
and private schools and a reproducible analysis pipeline; the four contributions
of \autoref{sec:intro} are delivered except for the survey results. The survey has
not been conducted, so RQ1--RQ3 remain open and no claim about teacher burnout is
made. For RQ1, the main outcome will be the difference in emotional exhaustion with
Cohen's $d$ and its bootstrap interval; for RQ2, the same outcome for each scale;
and for RQ3, the correlations and the regression of \autoref{tab:plan}.
The measurements, however, suggest that the pipeline is fast enough for repeated
use: in 120 runs without errors, the median latency stayed between 43.8 and
46.8\,ms for 35 to 1000 respondents.

Future work has three directions: (i) to collect the data and report the results
for RQ1--RQ3; (ii) to repeat the measurements on other machines with more runs,
so that tail percentiles become reliable; and (iii) to include more schools and
regions to test whether the findings generalise.
